\chapter{§3.3 平行与垂直}
\label{sec:3.3}


\prerequisite{§3.2}

\gradelevel{7 年级}


\section*{平行线的概念}


\begin{explore}


铁轨的两条钢轨永远不会交叉——无论它们向前延伸多远，始终保持相同的间距。这就是"平行"的直观含义。

\end{explore}

\begin{understand}


在同一平面内，\textbf{不相交}的两条直线叫做\textbf{平行线}。直线 $a$ 与直线 $b$ 平行，记作 $a \parallel b$ 。

\begin{center}
\begin{tikzpicture}[every node/.style={font=\small}]
  % Two parallel lines
  \draw[{Stealth[length=5pt]}-{Stealth[length=5pt]}, thick] (-0.3,1.5) -- (4.5,1.5) node[right] {$a$};
  \draw[{Stealth[length=5pt]}-{Stealth[length=5pt]}, thick] (-0.3,0)   -- (4.5,0)   node[right] {$b$};
  % Parallel tick marks
  \foreach \y in {0,1.5} {
    \draw (1.9,\y-0.12) -- (2.1,\y+0.12);
    \draw (2.1,\y-0.12) -- (2.3,\y+0.12);
  }
  \node[left] at (-0.3,0.75) {$a \parallel b$};
\end{tikzpicture}
\captionof{figure}{平行线 $a \parallel b$：同一平面内不相交的两条直线}
\label{fig:parallel-lines}
\end{center}


平行公理：经过直线外一点，有且只有一条直线与已知直线平行。

由此推论：如果两条直线都与第三条直线平行，那么这两条直线也互相平行。即：若 $a \parallel b$ ， $b \parallel c$ ，则 $a \parallel c$ 。

\end{understand}

\begin{keytakeaway}


\begin{itemize}
  \item 同一平面内，不相交的两条直线是平行线。
  \item 经过直线外一点，有且只有一条直线与已知直线平行。
  \item 平行具有传递性。
\end{itemize}


\end{keytakeaway}



\section*{平行线的判定}


\begin{explore}


工人砌墙时会用一根横线（水平线）和一把角尺来检查砖缝是否平行。他检查的其实就是某些角是否相等——这正是平行线判定的核心思想。

\end{explore}

\begin{understand}


当一条直线（称为\textbf{截线}）与两条直线相交时，会形成 $8$ 个角。这些角之间有特殊的名称和位置关系。

\begin{center}
\begin{tikzpicture}[every node/.style={font=\scriptsize}]
  % Lines a and b
  \draw[{Stealth[length=5pt]}-{Stealth[length=5pt]}, thick] (-0.5,2) -- (5,2) node[right,font=\small] {$a$};
  \draw[{Stealth[length=5pt]}-{Stealth[length=5pt]}, thick] (-0.5,0) -- (5,0) node[right,font=\small] {$b$};
  % Transversal l: from (0.5,3) to (3.5,-1)
  \draw[{Stealth[length=5pt]}-{Stealth[length=5pt]}, thick] (0.5,3) -- (3.5,-1) node[below,font=\small] {$l$};
  % Intersections: A at (1.25,2), B at (2.75,0)  [line: x=0.5+t*3, y=3-t*4 → y=2: t=0.25→x=1.25; y=0: t=0.75→x=2.75]
  \coordinate (A) at (1.25,2);
  \coordinate (B) at (2.75,0);
  \fill (A) circle (1.5pt) node[above right, font=\small] {$A$};
  \fill (B) circle (1.5pt) node[above right, font=\small] {$B$};
  % Angle labels at A (upper intersection)
  \node[mathred]   at (0.55,2.2)  {$\angle 1$};
  \node            at (1.8,2.2)   {$\angle 2$};
  \node            at (1.7,1.75)  {$\angle 3$};
  \node[mathblue]  at (0.6,1.75)  {$\angle 4$};
  % Angle labels at B (lower intersection)
  \node[mathblue]  at (2.1,0.2)   {$\angle 5$};
  \node            at (3.3,0.2)   {$\angle 6$};
  \node[mathred]   at (3.25,-0.25){$\angle 7$};
  \node            at (2.15,-0.25){$\angle 8$};
\end{tikzpicture}
\captionof{figure}{截线 $l$ 与直线 $a$、$b$ 相交形成 8 个角：同位角（如 $\angle 1$ 和 $\angle 5$），内错角（如 $\angle 4$ 和 $\angle 6$），同旁内角（如 $\angle 3$ 和 $\angle 5$）}
\label{fig:transversal-angles}
\end{center}


设截线 $l$ 分别交直线 $a$ 、 $b$ 于点 $A$ 、 $B$ 。标记八个角后，我们可以识别三类特殊的角对：

\begin{center}
\begin{tabular}{lll}
\toprule
角对类型 & 位置特征 & 示例 \\
\midrule
\textbf{同位角} & 在截线同侧，在两条被截线的同方位 & $\angle 1$ 与 $\angle 5$ \\
\textbf{内错角} & 在截线两侧，在两条被截线之间 & $\angle 3$ 与 $\angle 5$ \\
\textbf{同旁内角} & 在截线同侧，在两条被截线之间 & $\angle 3$ 与 $\angle 6$ \\
\bottomrule
\end{tabular}
\end{center}


\textbf{平行线的判定方法}（由角的关系判断两直线平行）：

\begin{enumerate}
  \item \textbf{同位角相等，两直线平行。}
  \item \textbf{内错角相等，两直线平行。}
  \item \textbf{同旁内角互补，两直线平行。}
\end{enumerate}


这三条判定方法的核心逻辑是一致的：只要其中一条成立，另外两条也必然成立。

\end{understand}

\begin{workedexamples}


\textbf{例 1.} 如图，直线 $l$ 截直线 $a$ 、 $b$ ， $\angle 1 = \angle 2$ 。试说明 $a \parallel b$ 。

\textbf{解：}
因为 $\angle 1 = \angle 2$ （已知），

又因为 $\angle 1$ 与 $\angle 2$ 是同位角，

所以 $a \parallel b$ （同位角相等，两直线平行）。

\textbf{例 2.} 如图， $\angle BAC = 70°$ ， $\angle ACD = 110°$ ，判断 $AB$ 与 $CD$ 的位置关系并说明理由。

\begin{center}
\begin{tikzpicture}[every node/.style={font=\small}]
  % Lines AB and CD with transversal AC
  \coordinate (A) at (0,2);
  \coordinate (B) at (3.5,2);
  \coordinate (C) at (0.8,0);
  \coordinate (D) at (4,0);
  \draw[{Stealth[length=5pt]}-{Stealth[length=5pt]}, thick] (-0.3,2) -- (3.8,2);
  \draw[{Stealth[length=5pt]}-{Stealth[length=5pt]}, thick] (0.2,0) -- (4.3,0);
  \draw[main line] (A) -- (C);
  \node[above left] at (A) {$A$};
  \node[above right] at (B) {$B$};
  \node[below left] at (C) {$C$};
  \node[below right] at (D) {$D$};
  % Angle marks
  \draw[angle arc, mathred, thick]
    ($(A)+(0:0.55)$) arc[start angle=0, end angle={-atan2(2,0.8)}, radius=0.55];
  \node[mathred, font=\scriptsize] at ($(A)+(333:0.85)$) {$70°$};
  \draw[angle arc, mathblue, thick]
    ($(C)+(90:0.55)$) arc[start angle=90, end angle={180-atan2(2,0.8)}, radius=0.55];
  \node[mathblue, font=\scriptsize] at ($(C)+(130:0.9)$) {$110°$};
\end{tikzpicture}
\captionof{figure}{同旁内角之和 $= 70° + 110° = 180°$，故 $AB \parallel CD$}
\label{fig:parallel-judge-ex2}
\end{center}


\textbf{解：}
因为 $\angle BAC + \angle ACD = 70° + 110° = 180°$ ，

又因为 $\angle BAC$ 与 $\angle ACD$ 是截线 $AC$ 所截的同旁内角，

所以 $AB \parallel CD$ （同旁内角互补，两直线平行）。

\textbf{例 3.} 如图， $\angle 1 = 45°$ ， $\angle 2 = 45°$ ，判断直线 $a$ 与直线 $b$ 的位置关系并说明理由。

\textbf{解：}
因为 $\angle 1 = \angle 2 = 45°$ （已知），

又因为 $\angle 1$ 与 $\angle 2$ 是内错角，

所以 $a \parallel b$ （内错角相等，两直线平行）。

\end{workedexamples}

\begin{keytakeaway}


\begin{itemize}
  \item 平行线有三种判定方法：同位角相等、内错角相等、同旁内角互补。
  \item 判定是"由角的关系推出平行"。
\end{itemize}


\end{keytakeaway}

\begin{practice}


\begin{enumerate}
  \item 如图，直线 $l$ 截直线 $a$ 、 $b$ ， $\angle 1 = 120°$ ， $\angle 2 = 60°$ 。直线 $a$ 与 $b$ 平行吗？说明理由。
  \item 如图， $\angle A = 50°$ ， $\angle C = 50°$ ， $\angle A$ 与 $\angle C$ 是截线 $AC$ 所截的内错角。 $AB$ 与 $CD$ 平行吗？说明理由。
\end{enumerate}


\end{practice}



\section*{平行线的性质}


\begin{explore}


判定告诉我们"什么条件下两直线平行"。现在反过来问：\textbf{已知两直线平行}，能得出什么结论？

\end{explore}

\begin{understand}


\textbf{平行线的性质}（由平行关系推出角的关系）：

\begin{enumerate}
  \item \textbf{两直线平行，同位角相等。}
  \item \textbf{两直线平行，内错角相等。}
  \item \textbf{两直线平行，同旁内角互补。}
\end{enumerate}


\end{understand}

\begin{quote}
\textbf{性质与判定的区别}：
- \textbf{判定}：角的关系 $\Rightarrow$ 平行（用来"证明平行"）
- \textbf{性质}：平行 $\Rightarrow$ 角的关系（用来"求角的度数"）

两者方向相反，不可混淆。
\end{quote}


\begin{workedexamples}


\textbf{例 1.} 如图， $a \parallel b$ ，截线 $l$ 与直线 $a$ 交于点 $A$ ，与直线 $b$ 交于点 $B$ 。 $\angle 1 = 65°$ ，求 $\angle 2$ 和 $\angle 3$ 。

\begin{center}
\begin{tikzpicture}[every node/.style={font=\small}]
  \draw[{Stealth[length=5pt]}-{Stealth[length=5pt]}, thick] (-0.3,2) -- (4.5,2) node[right] {$a$};
  \draw[{Stealth[length=5pt]}-{Stealth[length=5pt]}, thick] (-0.3,0) -- (4.5,0) node[right] {$b$};
  \draw[{Stealth[length=5pt]}-{Stealth[length=5pt]}, thick] (0.8,3) -- (3,-1) node[below] {$l$};
  % Intersections at A(1.5,2) and B(2.25,0)
  \coordinate (A) at (1.5,2);
  \coordinate (B) at (2.25,0);
  \fill (A) circle (2pt) node[above left] {$A$};
  \fill (B) circle (2pt) node[above left] {$B$};
  % Angle 1 at A (above-right of transversal), angle 2 at B (same side = corresponding)
  \node[mathred] at (1.95,2.25) {$\angle 1$};
  \node[mathblue] at (2.75,0.25) {$\angle 2$};
  \node[mathgreen] at (1.6,1.65) {$\angle 3$};
\end{tikzpicture}
\captionof{figure}{$a \parallel b$，$\angle 1 = 65°$；同位角 $\angle 2 = 65°$；同旁内角 $\angle 3 = 115°$}
\label{fig:parallel-property-ex1}
\end{center}


\textbf{解：}
因为 $a \parallel b$ （已知），

所以 $\angle 2 = \angle 1 = 65°$ （两直线平行，同位角相等）。

因为 $a \parallel b$ （已知），

所以 $\angle 1 + \angle 3 = 180°$ （两直线平行，同旁内角互补）。

所以 $\angle 3 = 180° - 65° = 115°$ 。

\textbf{例 2.} 如图， $AB \parallel CD$ ， $\angle ABE = 50°$ ， $\angle DCE = 30°$ 。求 $\angle BEC$ 。

\begin{center}
\begin{tikzpicture}[every node/.style={font=\small}]
  % AB and CD parallel lines
  \coordinate (A) at (0,2); \coordinate (B) at (4,2);
  \coordinate (C) at (0,0); \coordinate (D) at (4,0);
  \draw[{Stealth[length=5pt]}-{Stealth[length=5pt]},thick] (-0.3,2)--(4.3,2);
  \draw[{Stealth[length=5pt]}-{Stealth[length=5pt]},thick] (-0.3,0)--(4.3,0);
  \node[above left] at (0,2) {$A$}; \node[above right] at (4,2) {$B$};
  \node[below left] at (0,0) {$C$}; \node[below right] at (4,0) {$D$};
  % Point E between the lines
  \coordinate (E) at (2.5,1);
  \fill (E) circle (2pt) node[right] {$E$};
  \draw[main line] (B)--(E)--(C);
  % Angle marks
  \draw[angle arc,mathred,thick]
    ($(B)+(180:0.6)$) arc[start angle=180,end angle={180+atan2(-1,1.5)},radius=0.6];
  \node[mathred,font=\scriptsize] at ($(B)+(200:1.0)$) {$50°$};
  \draw[angle arc,mathblue,thick]
    ($(C)+(0:0.6)$) arc[start angle=0,end angle={atan2(1,2.5)},radius=0.6];
  \node[mathblue,font=\scriptsize] at ($(C)+(15:1.0)$) {$30°$};
\end{tikzpicture}
\captionof{figure}{$AB \parallel CD$，$\angle ABE=50°$，$\angle DCE=30°$，求 $\angle BEC$}
\label{fig:parallel-property-ex2}
\end{center}


\textbf{解：}
过点 $E$ 作 $EF \parallel AB$ 。

因为 $AB \parallel CD$ （已知）， $EF \parallel AB$ （作法），

所以 $EF \parallel CD$ （平行的传递性）。

因为 $EF \parallel AB$ ，

所以 $\angle BEF = \angle ABE = 50°$ （两直线平行，内错角相等）。

因为 $EF \parallel CD$ ，

所以 $\angle CEF = \angle DCE = 30°$ （两直线平行，内错角相等）。

所以 $\angle BEC = \angle BEF + \angle CEF = 50° + 30° = 80°$ 。

\textbf{例 3.} 如图， $AB \parallel CD$ ， $\angle A = 40°$ ， $\angle C = 30°$ ，求 $\angle AEC$ 。

\begin{center}
\begin{tikzpicture}[every node/.style={font=\small}]
  % AB (upper) and CD (lower) parallel
  \draw[{Stealth[length=5pt]}-{Stealth[length=5pt]},thick] (-0.3,2)--(4.3,2);
  \draw[{Stealth[length=5pt]}-{Stealth[length=5pt]},thick] (-0.3,0)--(4.3,0);
  \node[above] at (-0.1,2) {$A$}; \node[above right] at (4.3,2) {$B$};
  \node[below] at (-0.1,0) {$C$}; \node[below right] at (4.3,0) {$D$};
  % E is a point between the lines; AE and CE are drawn
  \coordinate (AE_start) at (0.5,2);
  \coordinate (CE_start) at (1.0,0);
  \coordinate (E) at (2.0,1.0);
  \fill (E) circle (2pt) node[right] {$E$};
  \draw[main line] (AE_start)--(E)--(CE_start);
  % Angle A = 40°
  \draw[angle arc,mathred,thick]
    ($(AE_start)+(0:0.5)$) arc[start angle=0,end angle={-atan2(1,1.5)},radius=0.5];
  \node[mathred,font=\scriptsize] at ($(AE_start)+(-15:0.85)$) {$40°$};
  % Angle C = 30°
  \draw[angle arc,mathblue,thick]
    ($(CE_start)+(0:0.5)$) arc[start angle=0,end angle={atan2(1,1.0)},radius=0.5];
  \node[mathblue,font=\scriptsize] at ($(CE_start)+(25:0.85)$) {$30°$};
\end{tikzpicture}
\captionof{figure}{$AB \parallel CD$，$\angle A=40°$，$\angle C=30°$，求 $\angle AEC$（答：$70°$）}
\label{fig:parallel-property-ex3}
\end{center}


\textbf{解：}
过点 $E$ 作 $EF \parallel AB$ ，则 $EF \parallel CD$ 。

因为 $EF \parallel AB$ ，所以 $\angle AEF = \angle A = 40°$ （内错角相等）。

因为 $EF \parallel CD$ ，所以 $\angle CEF = \angle C = 30°$ （内错角相等）。

所以 $\angle AEC = \angle AEF + \angle CEF = 40° + 30° = 70°$ 。

\end{workedexamples}

\begin{keytakeaway}


\begin{itemize}
  \item 两直线平行 $\Rightarrow$ 同位角相等、内错角相等、同旁内角互补。
  \item 遇到折线型角度问题，常用辅助线"过折点作平行线"。
\end{itemize}


\end{keytakeaway}

\begin{practice}


\begin{enumerate}
  \item $a \parallel b$ ，截线 $l$ 截出的 $\angle 1 = 72°$ ，求 $\angle 1$ 的内错角和同旁内角。
  \item 如图， $AB \parallel CD$ ， $\angle B = 65°$ ， $\angle D = 40°$ 。求 $\angle BED$ 。
  \item 如图， $AB \parallel CD$ ， $BE$ 平分 $\angle ABC$ ， $\angle BCD = 70°$ ，求 $\angle BEC$ 。
\end{enumerate}


\end{practice}



\section*{垂直的概念}


\begin{explore}


建筑工人用铅垂线来检查墙壁是否竖直——铅垂线是竖直向下的，而地面是水平的，它们之间的夹角恰好是 $90°$ 。当两条直线的夹角是 $90°$ 时，我们说它们\textbf{垂直}。

\end{explore}

\begin{understand}


当两条直线相交所成的角为 $90°$ （即直角）时，这两条直线\textbf{互相垂直}。其中一条直线叫做另一条直线的\textbf{垂线}，交点叫做\textbf{垂足}。

直线 $a$ 与直线 $b$ 垂直，记作 $a \perp b$ 。

\begin{center}
\begin{tikzpicture}[every node/.style={font=\small}]
  % Horizontal line a
  \draw[{Stealth[length=5pt]}-{Stealth[length=5pt]},thick] (-1.8,0)--(1.8,0) node[right] {$a$};
  % Vertical line b
  \draw[{Stealth[length=5pt]}-{Stealth[length=5pt]},thick] (0,-1.2)--(0,1.8) node[above] {$b$};
  % Right angle mark
  \draw (0.22,0)--(0.22,0.22)--(0,0.22);
  \fill (0,0) circle (2pt);
  \node[below right] at (0,0) {$O$};
  \node[mathred] at (0.6,0.5) {$90°$};
\end{tikzpicture}
\captionof{figure}{$a \perp b$：直线 $a$ 与直线 $b$ 垂直，垂足 $O$，交角均为直角}
\label{fig:perpendicular}
\end{center}


性质：\textbf{过一点有且只有一条直线与已知直线垂直。}

\end{understand}

\begin{quote}
注意和平行公理的对比：
- 过直线\textbf{外}一点，有且只有一条直线与已知直线\textbf{平行}。
- 过一点（直线上或直线外），有且只有一条直线与已知直线\textbf{垂直}。
\end{quote}


\begin{workedexamples}


\textbf{例 1.} 两条直线相交，如果其中一个角是直角，那么其他三个角分别是多少？

\textbf{解：}
设 $\angle 1 = 90°$ 。

$\angle 2 = 180° - 90° = 90°$ （邻补角）。

$\angle 3 = \angle 1 = 90°$ （对顶角）。

$\angle 4 = \angle 2 = 90°$ （对顶角）。

所以四个角都是 $90°$ 。这就是说：\textbf{两条直线垂直时，所形成的四个角都是直角。}

\end{workedexamples}

\begin{keytakeaway}


\begin{itemize}
  \item 两直线相交成 $90°$ 角，则互相垂直。
  \item 垂直时四个角都是直角。
  \item 过一点有且只有一条直线与已知直线垂直。
\end{itemize}


\end{keytakeaway}



\section*{垂线段最短与点到直线的距离}


\begin{explore}


站在河边想去对岸，怎样走最近？当然是走"垂直于河岸"的路线。这背后的原理就是\textbf{垂线段最短}。

\end{explore}

\begin{understand}


从直线外一点 $P$ 向直线 $l$ 引一条垂线，垂足为 $H$ 。再从 $P$ 向 $l$ 上的其他任意点 $A$ 连线段。

\begin{center}
\begin{tikzpicture}[every node/.style={font=\small}]
  % Line l
  \draw[{Stealth[length=5pt]}-{Stealth[length=5pt]},thick] (-0.5,0)--(5,0) node[right] {$l$};
  % Point P above line
  \coordinate (P) at (1.5,2.2);
  \fill (P) circle (2pt) node[above] {$P$};
  % Perpendicular foot H
  \coordinate (H) at (1.5,0);
  \fill (H) circle (2pt) node[below] {$H$};
  \draw[main line, mathblue] (P)--(H);
  \draw (1.72,0)--(1.72,0.22)--(1.5,0.22);
  \node[mathblue] at (1.1,1.1) {$PH$};
  % Other points A and B on line
  \coordinate (A) at (3.2,0);
  \coordinate (B) at (0.2,0);
  \fill (A) circle (2pt) node[below] {$A$};
  \fill (B) circle (2pt) node[below] {$B$};
  \draw[aux line] (P)--(A);
  \draw[aux line] (P)--(B);
  \node[mathred] at (2.7,1.3) {$PA > PH$};
\end{tikzpicture}
\captionof{figure}{垂线段最短：$PH \perp l$，对 $l$ 上任意 $A \neq H$，$PH < PA$；$PH$ 即点 $P$ 到直线 $l$ 的距离}
\label{fig:perpendicular-shortest}
\end{center}


结论：\textbf{从直线外一点到这条直线上各点的连线中，垂线段最短。}

即 $PH < PA$ （对于 $l$ 上任意不同于 $H$ 的点 $A$ ）。

\textbf{点到直线的距离}：从直线外一点到这条直线的垂线段的长度，叫做\textbf{点到直线的距离}。

\end{understand}

\begin{quote}
注意区分：
- "两点之间的距离"是一条线段的长度。
- "点到直线的距离"是垂线段的长度。
\end{quote}


\begin{workedexamples}


\textbf{例 1.} 如图， $P$ 是直线 $l$ 外一点， $PA$ 、 $PB$ 、 $PC$ 是从 $P$ 到 $l$ 上不同点的三条线段。其中 $PA \perp l$ ， $PA = 3$ cm， $PB = 4$ cm， $PC = 5$ cm。点 $P$ 到直线 $l$ 的距离是多少？

\textbf{解：}
因为 $PA \perp l$ ，所以 $PA$ 是从 $P$ 到 $l$ 的垂线段。

点 $P$ 到直线 $l$ 的距离就是垂线段 $PA$ 的长度，即 $3$ cm。

\textbf{例 2.} 如图， $a \parallel b$ ，点 $P$ 在直线 $a$ 上， $PH \perp b$ ，垂足为 $H$ ， $PH = 5$ cm。在直线 $a$ 上另取一点 $Q$ ， $QK \perp b$ ，垂足为 $K$ 。 $QK$ 一定也是 $5$ cm 吗？

\textbf{解：}
是的。因为 $a \parallel b$ ，两平行线之间的距离处处相等。无论在 $a$ 上取哪一点作 $b$ 的垂线段，其长度都相同。

这就是\textbf{平行线之间的距离}：两条平行线间的公垂线段的长度叫做两条平行线之间的距离。

\textbf{例 3.} $\triangle ABC$ 中， $AB = 5$ cm， $AC = 4$ cm， $BC = 6$ cm。从 $A$ 到 $BC$ 的距离指的是什么？

\textbf{解：}
从点 $A$ 向直线 $BC$ 作垂线，垂足为 $D$ ，则线段 $AD$ 的长度就是点 $A$ 到直线 $BC$ 的距离。 $AD$ 也叫做 $\triangle ABC$ 中 $BC$ 边上的\textbf{高}。

\end{workedexamples}

\begin{keytakeaway}


\begin{itemize}
  \item 从直线外一点到直线的所有连线中，垂线段最短。
  \item 点到直线的距离 = 垂线段的长度。
  \item 两平行线之间的距离处处相等。
  \item 三角形中，顶点到对边的垂线段就是"高"。
\end{itemize}


\end{keytakeaway}

\begin{practice}


\begin{enumerate}
  \item 点 $P$ 到直线 $l$ 的距离为 $4$ cm。在直线 $l$ 上取一点 $A$ ，使得 $PA = 5$ cm。这样的点 $A$ 有几个？
  \item 两条平行线之间的距离为 $8$ cm。在一条平行线上取一点 $P$ ，在另一条上取一点 $Q$ 。 $PQ$ 的最小值是多少？
  \item 如图， $\angle AOB = 30°$ ， $P$ 是 $\angle AOB$ 内部一点， $P$ 到 $OA$ 的距离是 $3$ cm。求 $OP$ 的长。
\end{enumerate}


\end{practice}



\begin{answer}


\begin{enumerate}
  \item $\angle 1 + \angle 2 = 120° + 60° = 180°$ 。因为 $\angle 1$ 与 $\angle 2$ 是同旁内角且互补，所以 $a \parallel b$ 。
\end{enumerate}


\begin{enumerate}
  \item 因为 $\angle A = \angle C = 50°$ ，且 $\angle A$ 与 $\angle C$ 是内错角，内错角相等，所以 $AB \parallel CD$ 。
\end{enumerate}


\begin{enumerate}
  \item 内错角 $= 72°$ （内错角相等）。同旁内角 $= 180° - 72° = 108°$ （同旁内角互补）。
\end{enumerate}


\begin{enumerate}
  \item 过点 $E$ 作 $EF \parallel AB$ ，则 $EF \parallel CD$ 。 $\angle BEF = \angle B = 65°$ （内错角）， $\angle DEF = \angle D = 40°$ （内错角）。 $\angle BED = \angle BEF + \angle DEF = 65° + 40° = 105°$ 。
\end{enumerate}


\begin{enumerate}
  \item 因为 $AB \parallel CD$ ，所以 $\angle ABC + \angle BCD = 180°$ （同旁内角互补）， $\angle ABC = 180° - 70° = 110°$ 。因为 $BE$ 平分 $\angle ABC$ ，所以 $\angle ABE = 55°$ 。过 $E$ 作 $EF \parallel AB$ ，则 $EF \parallel CD$ ， $\angle BEF = \angle ABE = 55°$ （内错角）， $\angle CEF = \angle DCE = 70°$ （内错角）。 $\angle BEC = \angle BEF + \angle CEF = 55° + 70° = 125°$ 。
\end{enumerate}


\begin{enumerate}
  \item 设垂足为 $H$ ， $PH = 4$ cm。在直线 $l$ 上，满足 $PA = 5$ cm 的点 $A$ 满足 $HA = \sqrt{PA^2 - PH^2} = \sqrt{25 - 16} = 3$ cm（由勾股定理，此处可直观理解为以 $P$ 为圆心、 $5$ 为半径的圆与直线 $l$ 的交点）。 $H$ 两侧各有一个，所以共 $2$ 个。
\end{enumerate}


\begin{enumerate}
  \item $PQ$ 的最小值就是两平行线之间的距离，即 $PQ_{\min} = 8$ cm（此时 $PQ \perp$ 两平行线）。
\end{enumerate}


\begin{enumerate}
  \item 设 $P$ 到 $OA$ 的垂足为 $H$ ，则 $PH = 3$ cm， $\angle POH = 30°$ （或更小，但这里 $PH \perp OA$ ）。在直角 $\triangle OPH$ 中， $\sin 30° = \dfrac{PH}{OP}$ ，即 $\dfrac{1}{2} = \dfrac{3}{OP}$ ， $OP = 6$ cm。（此题用到了直角三角形中角与边的关系，详见 §3.10。在当前阶段可记住： $30°$ 角的对边等于斜边的一半，所以 $PH = \dfrac{1}{2} OP$ ， $OP = 2 \times 3 = 6$ cm。）
\end{enumerate}


\end{answer}
